% =============================================================================
% CAPÍTULO 11: Validação Cruzada com K-Fold no SUS-Twin
% Volume 2 — Métricas de validação clínica, K-Fold temporal,
%            TRIPOD-AI e interpretação de resultados
% =============================================================================
% V2 Standalone — sem \input{}
% Seções:
%   11.1 - Introdução
%   11.2 - Preparação de dados clínicos com DATASUS e PySUS
%   11.3 - K-Fold tradicional vs. validação temporal
%   11.4 - Métricas TRIPOD-AI: sensibilidade, especificidade, AUC, R²
%   11.5 - Implementação prática com Periomod e Scikit-learn
%   11.6 - Exercícios propostos
% =============================================================================

\destaque{Nível-alvo N2 — Pesquisa Reprodutível: implemente validação cruzada
temporal com K-Fold e métricas TRIPOD-AI seguindo o Framework SUS-Twin.}

\section{Introdução}
\label{ch:ia}

A validação de modelos preditivos em odontologia digital não é um
complemento opcional: é o requisito ético e científico que separa uma
inferência clínica confiável de uma correlação espúria. No contexto do
Framework SUS-Twin, um gêmeo digital periodontal que não passa por
validação cruzada rigorosa pode produzir predições otimistas,
superajustadas (\textit{overfitting}\footnote{Overfitting (superajuste)
ocorre quando o modelo decora os dados de treino em vez de aprender
padrões generalizáveis. Um modelo superajustado apresenta desempenho
excelente nos dados que já viu, mas falha catastrophicamente ao
encontrar novos pacientes.}) aos dados de treino e, portanto,
clinicamente perigosas.

A validação cruzada com K-Fold\footnote{K-Fold é uma técnica
estatística que divide o conjunto de dados em $k$ partes (\"folds\").
O modelo é treinado em $k-1$ partes e testado na parte restante,
repetindo o processo $k$ vezes. Cada observação é usada para teste
exatamente uma vez, gerando uma estimativa robusta do desempenho real
do modelo.} é o padrão-ouro para estimar o erro de generalização de
modelos de aprendizado de máquina em saúde. Quando aplicada a séries
temporais de dados periodontais --- como a progressão da perda de
inserção clínica (PIC) ao longo de consultas sucessivas no SUS --- a
versão \textbf{temporal} do K-Fold respeita a ordem cronológica dos
eventos, evitando que o modelo \"aprenda o futuro\" para predizer o
passado.

Este capítulo apresenta, em ordem sequencial: (1) a preparação de
dados clínicos reais do DATASUS usando a biblioteca PySUS; (2) a
distinção conceitual e prática entre K-Fold tradicional e temporal;
(3) as métricas do guideline TRIPOD-AI\footnote{TRIPOD-AI
(Transparent Reporting of a multivariable prediction model for
Individual Prognosis Or Diagnosis — Artificial Intelligence) é a
diretriz internacional para relato transparente de modelos de
predição baseados em IA. Define métricas obrigatórias como
sensibilidade, especificidade, AUC e R², além de exigir validação
interna e externa.} com tabela de fórmulas e interpretação clínica;
(4) a implementação prática de um pipeline de validação cruzada
usando a biblioteca \texttt{Periomod}\footnote{Periomod é uma
biblioteca Python especializada em modelagem preditiva periodontal.
Implementa algoritmos de sobrevivência, regressão longitudinal e
validação cruzada especificamente calibrados para métricas
odontológicas, como perda de inserção clínica (PIC) e sangramento à
sondagem (SS).} integrada ao \texttt{scikit-learn}; e (5) três
exercícios progressivos que consolidam o aprendizado.

\section{Preparação de Dados Clínicos: DATASUS e PySUS}

O primeiro passo de qualquer pipeline de validação cruzada é obter
dados clínicos representativos e estruturados. No Brasil, a principal
fonte de dados epidemiológicos periodontais é o
DATASUS\footnote{DATASUS é o departamento de informática do Sistema
Único de Saúde que disponibiliza publicamente os registros de todos
os atendimentos ambulatoriais e hospitalares do país, incluindo
procedimentos odontológicos.}, acessado programaticamente pela
biblioteca PySUS\footnote{A PySUS é uma interface Python que baixa e
processa arquivos dos sistemas SIA (Ambulatorial), SIH (Hospitalar) e
SIM (Mortalidade) do DATASUS, convertendo formatos proprietários
(\texttt{.dbc}) em DataFrames pandas prontos para análise.}.

O código abaixo demonstra o download de dados de procedimentos
periodontais do estado de São Paulo, a filtragem por códigos
específicos de periodontia e a estruturação em formato tabular com a
biblioteca \texttt{pandas}\footnote{pandas é a biblioteca padrão do
Python para manipulação de dados tabulares. Oferece estruturas como
DataFrame (tabela com linhas e colunas) e Series (vetor), além de
operações de filtro, agrupamento e estatística descritiva.}:

\begin{lstlisting}[language=Python, caption={Download e estruturação de
dados periodontais do DATASUS com PySUS},
label={lst:pysus_perio}]
from pysus import SIA
import pandas as pd
import numpy as np

def baixar_dados_periodontais(uf: str, ano: int, mes: int) -> pd.DataFrame:
    """
    Baixa a producao ambulatorial odontologica de um estado
    e filtra apenas procedimentos periodontais.
    """
    sia = SIA().load(group='PA', state=uf, year=ano, month=mes)
    
    # Filtro por codigos de periodontia (PROCID)
    # Fonte: tabela unificada do SUS - CBOD
    codigos_perio = [
        '0301010012',  # Raspagem alveolar
        '0301010039',  # Curetagem periodontal
        '0301010055',  # Gengivectomia
        '0301010071',  # Aumento de coroa clinica
        '0301010098',  # Retalho periodontal
    ]
    
    perio = sia[sia['PROCID'].isin(codigos_perio)].copy()
    
    # Engenharia de atributos clinicos
    perio['mes_ano'] = pd.to_datetime(
        perio['ANO_CMPT'].astype(str) + '-' +
        perio['MES_CMPT'].astype(str) + '-01'
    )
    
    # Agregacao por estabelecimento e mes
    agregado = perio.groupby(
        ['mes_ano', 'CO_UF', 'CO_MUN', 'CO_CNES']
    ).agg(
        total_procedimentos=('PROCID', 'count'),
        valor_total=('VAL_APR', 'sum'),
        tipos_distintos=('PROCID', 'nunique')
    ).reset_index()
    
    return agregado

# Exemplo: dados de SP em janeiro de 2025
df_perio = baixar_dados_periodontais('SP', 2025, 1)
print(f"Registros: {len(df_perio)}")
print(f"Municipios: {df_perio['CO_MUN'].nunique()}")
print(f"Valor total: R$ {df_perio['valor_total'].sum():,.2f}")
\end{lstlisting}

O DataFrame resultante contém, para cada estabelecimento de saúde no
estado selecionado, o volume mensal de procedimentos periodontais, o
valor desembolsado pelo SUS e a diversidade de tipos de
procedimentos. Esta estrutura tabular é a matéria-prima para os
modelos preditivos que validaremos nas seções seguintes.

\section{K-Fold Tradicional vs. Validação Temporal}

\subsection{K-Fold Tradicional}

A validação cruzada K-Fold\footnote{K-Fold é uma técnica de
reamostragem que particiona o conjunto de dados em $k$ subconjuntos
mutuamente exclusivos de tamanho aproximadamente igual. O modelo é
treinado em $k-1$ folds e testado no fold restante, rotacionando o
fold de teste até que todos os $k$ folds tenham servido como teste. O
desempenho final é a média das $k$ iterações.}
particiona os dados em $k$ dobras mutuamente exclusivas. Para cada
iteração $i = 1, \dots, k$, o modelo é treinado nas $k-1$ dobras
restantes e testado na dobra $i$. A estimativa final de desempenho é
a média aritmética das $k$ métricas observadas:

\begin{equation}
    \text{Desempenho}_{\text{K-Fold}} = \frac{1}{k} \sum_{i=1}^{k} \mathcal{M}(y_i, \hat{y}_i)
    \label{eq:kfold_media}
\end{equation}

onde $\mathcal{M}$ é a métrica de interesse (acurácia, AUC, R²) e
$y_i$, $\hat{y}_i$ são os valores reais e preditos do fold $i$.

A principal limitação do K-Fold tradicional em dados clínicos
longitudinais é a \textbf{contaminação temporal}: observações de um
mesmo paciente em meses distintos podem ser alocadas simultaneamente
nos conjuntos de treino e teste, fazendo com que o modelo \"veja\" o
futuro do paciente durante o treinamento. Isso infla
artificialmente as métricas e esconde a verdadeira capacidade
preditiva do modelo.

\subsection{K-Fold Temporal}

O K-Fold temporal resolve esse problema respeitando a ordem
cronológica das observações. Em vez de sortear as dobras
aleatoriamente, a partição é feita por janelas de tempo: o modelo é
treinado com dados de meses $1, \dots, t$ e testado nos meses
$t+1, \dots, T$.

O código a seguir implementa ambas as estratégias usando
\texttt{scikit-learn}\footnote{scikit-learn (sklearn) é a biblioteca
de aprendizado de máquina mais utilizada no ecossistema Python.
Oferece implementações otimizadas de validação cruzada, métricas de
desempenho, pré-processamento e dezenas de algoritmos de
classificação e regressão.}:

\begin{lstlisting}[language=Python, caption={Comparação entre
K-Fold tradicional e temporal com scikit-learn},
label={lst:kfold_comparison}]
from sklearn.model_selection import (
    KFold, TimeSeriesSplit, cross_val_score
)
from sklearn.ensemble import RandomForestRegressor
from sklearn.metrics import mean_squared_error, r2_score
import numpy as np
import pandas as pd

def comparar_kfold(X: pd.DataFrame, y: pd.Series, n_splits: int = 5):
    """
    Compara K-Fold tradicional vs. temporal em um dataset
    de serie temporal periodontal.
    """
    # --- K-Fold tradicional (embaralhado) ---
    kf_trad = KFold(
        n_splits=n_splits,
        shuffle=True,
        random_state=42
    )
    
    # --- K-Fold temporal (preserva ordem) ---
    kf_temp = TimeSeriesSplit(n_splits=n_splits)
    
    modelo = RandomForestRegressor(
        n_estimators=100,
        max_depth=10,
        random_state=42
    )
    
    # Avaliacao com K-Fold tradicional
    scores_trad = cross_val_score(
        modelo, X, y,
        cv=kf_trad,
        scoring='r2'
    )
    
    # Avaliacao com K-Fold temporal
    scores_temp = cross_val_score(
        modelo, X, y,
        cv=kf_temp,
        scoring='r2'
    )
    
    print("=== Comparacao K-Fold ===")
    print(f"Tradicional: R² = {scores_trad.mean():.3f} "
          f"(+/- {scores_trad.std() * 2:.3f})")
    print(f"Temporal:    R² = {scores_temp.mean():.3f} "
          f"(+/- {scores_temp.std() * 2:.3f})")
    print(f"Diferenca (tendencioso - realista): "
          f"{(scores_trad.mean() - scores_temp.mean()):.3f}")
    
    return scores_trad, scores_temp
\end{lstlisting}

É esperado que o K-Fold tradicional produza métricas
sistematicamente superiores ao temporal, especialmente em séries com
tendência temporal (por exemplo, aumento da demanda por periodontia
ao longo dos anos). A diferença entre os dois scores é uma estimativa
do viés de contaminação temporal.

\section{Métricas TRIPOD-AI: Sensibilidade, Especificidade, AUC e
\texorpdfstring{R²}{R2}}

O guideline TRIPOD-AI estabelece quais métricas devem ser reportadas
para garantir transparência e reprodutibilidade em modelos preditivos
baseados em IA. Para problemas de classificação (ex.: risco alto vs.
baixo de progressão periodontal), as métricas centrais são
sensibilidade, especificidade e AUC\footnote{AUC (Area Under the ROC
Curve) mede a capacidade do modelo de distinguir entre classes
positivas e negativas em todos os limiares de decisão possíveis.
Varia de 0,5 (classificação aleatória) a 1,0 (classificação perfeita).
Valores acima de 0,8 são considerados clinicamente úteis.}. Para
problemas de regressão (ex.: predição da perda de inserção clínica em
milímetros), o R² é a métrica principal.

A \autoref{tab:tripod_metrics} consolida as cinco métricas
fundamentais com suas fórmulas, interpretações clínicas e critérios
de aceitação:

\begin{table}[htb]
\centering
\caption{Métricas TRIPOD-AI para validação de modelos periodontais.}
\label{tab:tripod_metrics}
\begin{adjustbox}{max width=\textwidth}
\begin{tabular}{lcl}
\toprule
\textbf{Métrica} & \textbf{Fórmula} & \textbf{Interpretação Clínica} \\
\midrule
Sensibilidade\footnote{Sensibilidade (também chamada de recall ou
taxa de verdadeiros positivos) é a proporção de pacientes com a
condição (ex.: periodontite estágio III) que o modelo consegue
identificar corretamente. Uma sensibilidade de 0,90 significa que
90\% dos casos reais foram detectados.}
& $\displaystyle \frac{VP}{VP + FN}$
& Percentual de casos positivos corretamente identificados.
\, Aceitável: $> 0{,}80$. \\
Especificidade\footnote{Especificidade é a proporção de pacientes
saudáveis que o modelo classifica corretamente como negativos. Uma
especificidade de 0,85 significa que 85\% dos pacientes sem a
condição foram corretamente liberados.}
& $\displaystyle \frac{VN}{VN + FP}$
& Percentual de casos negativos corretamente identificados.
\, Aceitável: $> 0{,}80$. \\
AUC-ROC & Integral da curva ROC &
Capacidade discriminativa global. Independe do limiar.
\, Aceitável: $> 0{,}80$. \\
Acurácia & $\displaystyle \frac{VP + VN}{VP + VN + FP + FN}$ &
Percentual total de acertos. Útil quando classes são balanceadas.
\, Aceitável: $> 0{,}85$. \\
R²\footnote{O coeficiente de determinação R² mede a proporção da
variância total da variável resposta que é explicada pelo modelo.
Varia de $-\infty$ (modelo pior que a média) a 1 (modelo perfeito).
Valores acima de 0,5 indicam moderada capacidade preditiva. Em
estudos periodontais, R² $> 0{,}3$ já é considerado relevante.}
& $\displaystyle 1 - \frac{\sum (y_i - \hat{y}_i)^2}{\sum (y_i - \bar{y})^2}$ &
Proporção da variância explicada. Para regressão de PIC.
\, Aceitável: $> 0{,}30$. \\
\bottomrule
\end{tabular}
\end{adjustbox}
\end{table}

O código abaixo calcula todas as métricas TRIPOD-AI para um modelo de
classificação de risco periodontal, incluindo a matriz de confusão
completa:

\begin{lstlisting}[language=Python, caption={Cálculo completo das
métricas TRIPOD-AI para classificação periodontal},
label={lst:tripod_metrics_code}]
from sklearn.metrics import (
    confusion_matrix, roc_auc_score,
    recall_score, precision_score,
    accuracy_score, r2_score,
    classification_report
)
import numpy as np

def relatorio_tripod_ai(y_true: np.ndarray, y_pred: np.ndarray,
                        y_prob: np.ndarray = None,
                        nome_modelo: str = "Modelo"):
    """
    Gera relatorio completo TRIPOD-AI com todas as metricas
    obrigatorias para validacao de modelos periodontais.
    """
    VP, FN, FP, VN = confusion_matrix(y_true, y_pred).ravel()
    
    metricas = {
        'Sensibilidade (Recall)': recall_score(y_true, y_pred),
        'Especificidade': VN / (VN + FP),
        'Acurácia': accuracy_score(y_true, y_pred),
        'Precisão': precision_score(y_true, y_pred),
    }
    
    if y_prob is not None:
        metricas['AUC-ROC'] = roc_auc_score(y_true, y_prob)
    
    print(f"\n=== RELATORIO TRIPOD-AI: {nome_modelo} ===")
    print(f"{'Metrica':<25} {'Valor':<10} {'Interpretacao'}")
    print("-" * 65)
    for nome, valor in metricas.items():
        nivel = ("EXCELENTE" if valor >= 0.90 else
                 "ACEITAVEL" if valor >= 0.80 else
                 "MODERADO"  if valor >= 0.60 else
                 "INSUFICIENTE")
        print(f"{nome:<25} {valor:.3f}{'':>5} {nivel}")
    
    print(f"\nMatriz de Confusao:")
    print(f"  VP={VP}  FN={FN}")
    print(f"  FP={FP}  VN={VN}")
    
    return metricas
\end{lstlisting}

\section{Implementação Prática com Periomod e Scikit-learn}

A biblioteca \texttt{Periomod} oferece uma camada de abstração
especializada para modelagem periodontal, integrando-se
transparentemente ao ecossistema \texttt{scikit-learn}. O pipeline
completo de validação cruzada temporal com Periomod é apresentado no
código a seguir:

\begin{lstlisting}[language=Python, caption={Pipeline de validação
cruzada temporal com Periomod},
label={lst:periomod_kfold}]
import periomod as pm
from periomod.validation import TemporalValidator
from periomod.metrics import periodontal_risk_score
from sklearn.ensemble import GradientBoostingClassifier
import pandas as pd
import numpy as np

def pipeline_validacao_temporal(
    df: pd.DataFrame,
    n_splits: int = 5,
    paciente_col: str = 'id_paciente',
    tempo_col: str = 'mes_ano',
    alvo_col: str = 'risco_progressao'
):
    """
    Pipeline completo de validacao cruzada temporal
    com Periomod para predicao de risco periodontal.
    """
    # 1. Instancia o validador temporal do Periomod
    validador = TemporalValidator(
        n_splits=n_splits,
        patient_col=paciente_col,
        time_col=tempo_col,
        gap_dias=90  # minimo 90 dias entre treino e teste
    )
    
    # 2. Prepara features e alvo
    feature_cols = [
        'idade', 'sexo', 'fumo_diario',
        'diabetes', 'pic_basal_mm', 'ss_basal_pct',
        'num_consultas_12m', 'raspagem_12m',
        'uso_fio_dental', 'renda_familiar'
    ]
    X = df[feature_cols].values
    y = df[alvo_col].values
    
    # 3. Configura o classificador
    modelo = GradientBoostingClassifier(
        n_estimators=200,
        max_depth=4,
        learning_rate=0.05,
        random_state=42
    )
    
    # 4. Executa validacao temporal
    resultados = validador.validate(
        modelo, X, y,
        return_models=False
    )
    
    # 5. Exibe resultados agregados
    print(f"\n=== VALIDACAO TEMPORAL PERIOMOD ===")
    print(f"Folds: {n_splits}")
    print(f"AUC medio:   {resultados['auc_mean']:.3f} "
          f"(+/- {resultados['auc_std']:.3f})")
    print(f"Acuracia:    {resultados['accuracy_mean']:.3f}")
    print(f"Risco medio: {resultados['risk_score_mean']:.3f}")
    
    # 6. Score composto periodontal (TRIPOD-AI)
    score_final = periodontal_risk_score(
        y_true=resultados['y_true_all'],
        y_pred=resultados['y_pred_all'],
        y_prob=resultados['y_prob_all']
    )
    print(f"Score Periodontal TRIPOD-AI: {score_final:.1f}/100")
    
    return resultados, score_final
\end{lstlisting}

O \texttt{TemporalValidator} do Periomod implementa internamente:
\begin{itemize}
\item Segmentação dos pacientes por janelas cronológicas sem
sobreposição.
\item Garantia de que o mesmo paciente não aparece simultaneamente em
treino e teste (\textit{patient-wise split}\footnote{O
patient-wise split (divisão por paciente) assegura que todas as
observações de um mesmo paciente estão exclusivamente no conjunto de
treino ou no conjunto de teste. Isso evita vazamento de informação
entre consultas do mesmo indivíduo.}).
\item Espaçamento mínimo (\textit{gap}) entre o último ponto de
treino e o primeiro ponto de teste para simular o horizonte clínico
real.
\end{itemize}

\section{Exercícios Propostos}

Os exercícios a seguir consolidam os conceitos apresentados neste
capítulo e seguem a progressão Graduação \(\to\) Mestrado \(\to\)
Doutorado quanto à profundidade investigativa.

\begin{exercicio}[Graduação — Preparação de Dataset Periodontal]
\label{ex:preparar_dataset}
Utilizando a biblioteca PySUS, baixe os dados de procedimentos
periodontais do estado de Minas Gerais (UF = \texttt{MG}) para o ano
de 2024. Filtre exclusivamente os códigos \texttt{0301010012}
(raspagem alveolar) e \texttt{0301010039} (curetagem periodontal).
Agregue os dados por município (\texttt{CO\_MUN}) e calcule:
\begin{enumerate}
\item O número total de procedimentos realizados em cada município.
\item O valor médio pago por procedimento (coluna \texttt{VAL\_APR}).
\item O município com maior volume de raspagens alveolares.
\end{enumerate}
\textbf{Dica:} Consulte o \autoref{lst:pysus_perio} como ponto de
partida. Substitua a UF e o ano pelos parâmetros solicitados.
\end{exercicio}

\begin{exercicio}[Mestrado — Implementação de K-Fold Temporal]
\label{ex:kfold_temporal}
Com base no dataset preparado no \autoref{ex:preparar_dataset},
implemente um pipeline de validação cruzada temporal com as seguintes
especificações:
\begin{enumerate}
\item Use \texttt{TimeSeriesSplit} do scikit-learn com
\(k = 4\) splits.
\item Treine um \texttt{RandomForestClassifier} com 150 árvores e
profundidade máxima 8.
\item Reporte as métricas TRIPOD-AI (sensibilidade, especificidade,
AUC) para cada fold e a média entre folds.
\item Compare os resultados com um K-Fold tradicional
(\texttt{KFold}, shuffle=True) e explique por que as métricas diferem.
\end{enumerate}
\textbf{Material de apoio:} \autoref{lst:kfold_comparison} e
\autoref{lst:tripod_metrics_code}.
\end{exercicio}

\begin{exercicio}[Doutorado — Interpretação Clínica da Matriz de
Confusão]
\label{ex:matriz_confusao}
A \autoref{tab:matriz_exemplo} apresenta a matriz de confusão de um
modelo de triagem periodontal testado em 500 pacientes do SUS:

\begin{table}[htb]
\centering
\caption{Matriz de confusão hipotética de um modelo de triagem
periodontal.}
\label{tab:matriz_exemplo}
\begin{adjustbox}{max width=0.6\textwidth}
\begin{tabular}{lcc}
\toprule
 & \textbf{Predito: Positivo} & \textbf{Predito: Negativo} \\
\midrule
\textbf{Real: Positivo} & 180 & 20 \\
\textbf{Real: Negativo} & 45  & 255 \\
\bottomrule
\end{tabular}
\end{adjustbox}
\end{table}

Com base na matriz:
\begin{enumerate}
\item Calcule manualmente: sensibilidade, especificidade, acurácia,
precisão e \textit{F1-score}\footnote{O F1-score é a média harmônica
entre precisão e sensibilidade: $F_1 = 2 \cdot
\frac{\text{precisão} \cdot \text{sensibilidade}}{\text{precisão} +
\text{sensibilidade}}$. É especialmente útil quando as classes são
desbalanceadas.}.
\item Interprete clinicamente cada métrica: o modelo é seguro para
uso como ferramenta de triagem no SUS? Justifique.
\item Qual métrica é mais crítica em um cenário de triagem
populacional: sensibilidade ou especificidade? Por quê?
\item Refaça os cálculos usando o código do
\autoref{lst:tripod_metrics_code} e compare com seus resultados
manuais.
\end{enumerate}
\end{exercicio}

\vfill

\begin{highlight}
\centering\small
\textbf{Badge do Capítulo}\\[4pt]
Nível-alvo: N2 --- Pesquisa Reprodutível\\
Habilidades: K-Fold temporal, métricas TRIPOD-AI,\\
\hspace*{4.5em}Periomod, interpretação clínica\\
Pré-requisitos: Python 3.11+, scikit-learn, Periomod\\
Comando de verificação:\\
\verb|tdd_validate.py --chapter 11|
\end{highlight}
